(1) 在 $z$ 基底中选择相位惯例，使 $|\boldsymbol n,+\rangle=\cos(\theta/2)|+z\rangle+\sin(\theta/2)|-z\rangle$。由投影规则
\[
P(n+)=|\langle\boldsymbol n,+|+z\rangle|^2=\cos^2(\theta/2).
\]
测得正值后的归一化状态就是 $|\boldsymbol n,+\rangle$。

(2) 演化算符为 $U(t)=e^{i\omega t\sigma_z/2}$。$|+x\rangle=(|+z\rangle+|-z\rangle)/\sqrt2$，于是
\[
P(x+\mid n+)=|\langle+x|U(t)|\boldsymbol n,+\rangle|^2
=\tfrac12[1+\sin\theta\cos(\omega t)].
\]
相对相位为 $e^{-i\omega t}$，其符号不影响余弦，但会影响 $y$ 方向的自旋分量。

(3) 忽略首次测量结果（非选择性测量）时，密度算符为
\[
\rho'=\Pi_{n+}|+z\rangle\langle+z|\Pi_{n+}
+\Pi_{n-}|+z\rangle\langle+z|\Pi_{n-}
=\tfrac12\bigl[I+\cos\theta\,(\boldsymbol n\cdot\boldsymbol\sigma)\bigr].
\]
绕 $z$ 轴进动后其 Bloch 向量的 $x$ 分量为 $\cos\theta\sin\theta\cos\omega t$，因此
\[
P_{\mathrm{nonselect}}(x+)=\tfrac12[1+\cos\theta\sin\theta\cos(\omega t)].
\]
若从未执行首次测量，初态始终为 $|+z\rangle$（仅得整体相位），故 $P_{\mathrm{no\ measure}}(x+)=1/2$。

(4) 当 $\theta=\pi/2$，首次测得 $n+$ 的概率为 $1/2$；选中 $n+$ 后的末次 $x+$ 概率为 $(1+\cos\omega t)/2$，未选择结果则为 $1/2$，完全不做第一次测量也为 $1/2$。虽然此特殊角下两类无条件概率巧合相同，但条件概率不同；一般倾角下非选择性结果与未测量结果并不相同。因此丢弃经典记录不等于撤销量子测量。
